\section{Introduction}
\label{sec:intro}

Sampling from a Boltzmann density $\pi(x)\propto\exp(-\mathcal{E}(x))$ when only the energy $\mathcal{E}$ and its gradient are available is a core problem in Bayesian inference, statistical physics and molecular simulation~\citep{noe2019boltzmann}.
Diffusion-based neural samplers have become a leading approach~\citep{zhang2022path,vargas2023denoising,berner2024optimal,akhound2024iterated}.
Among them, \emph{energy-matching} methods are attractive because they are simulation-free in their inner loop.
iDEM~\citep{akhound2024iterated} and its successors~\citep{bnem2024,iefm2024,cvsi2025,iwsm2025} regress a score network onto a Monte Carlo estimate of the noised score $\nabla\log p_t$, computed by self-normalized importance sampling (SNIS), at points obtained by noising samples from a replay buffer of model samples.

This paper starts from an observation about the \emph{population optimum} of these objectives. The SNIS target depends only on the noised point and fresh proposals, never on the buffer sample that produced it. So the loss is minimized by the same field $b_K=\mathbb{E}[\widehat S_K]$ \emph{whatever the buffer contains}, and $b_K\neq\nabla\log p_t$ for finite $K$.
Even a perfect sampler is therefore not a fixed point of training. The model is pulled towards a biased field, and the bias is largest at the moderate and high noise levels that set the shape and weights of the modes. More Monte Carlo samples shrink this bias only slowly, because the effective sample size of SNIS collapses exponentially with dimension and noise level.
This inconsistency has gone unmeasured on the benchmarks where it matters most, because there the noised score is unknown. We show how to measure it from reference samples alone, find it to be large on the standard molecular benchmarks, and remove it.

We point out that every energy-matching training pair already contains the missing ingredient: the clean buffer sample $x_0$ from which $x_t$ was generated. If the buffer were exact, $x_0$ would be an exact draw from the denoising posterior $\pi_t(\cdot\mid x_t)$. \textbf{Anchored Energy Matching (AEM)} adds $x_0$ as one extra particle to the SNIS estimate. This is the conditional importance sampling construction behind conditional SMC and i-SIR~\citep{andrieu2010particle,andrieu2018uniform}, here used to define a regression target instead of a Markov chain. The change is one line of code and one extra energy evaluation. Our contributions are as follows.

\begin{itemize}
\item \textbf{Diagnosis} (\cref{cor:fixed}, \cref{prop:diag}). An energy-matching objective whose target is a finite-$K$ self-normalized estimate built from $x_t$ and fresh proposals alone, as in iDEM, has a buffer-independent population minimizer $b_K\neq\nabla\log p_t$, so the target distribution is not a fixed point of training. We test iDEM; successors that keep a self-normalized target (e.g.\ iEFM, CVSI) share the structure, but we do not test them. We give two checks that detect this from samples of $\pi$ alone, without the noised score: a $z$-test of conditional unbiasedness for any training target, and a held-out denoising loss that ranks trained models by their score error. Both apply to targets, such as molecular systems, where the score is unknown.
\item \textbf{Evidence on standard benchmarks} (\cref{sec:experiments}). The test rejects iDEM's target at moderate and high noise on DW-4 and LJ-13 ($|z|$ up to $39$ and $115$); on DW-4 ten times more proposals do not make it pass. Trained on a buffer of \emph{exact} samples, iDEM learns a score with relative error $25$--$35$ at $\sigma_t\in\{1,2\}$ on a 32-dimensional many-well target (AEM: $0.07$--$0.14$), and this error matches the directly measured bias of its target. On DW-4 it learns a worse mean score than AEM at every noise level.
\item \textbf{A one-line fix with guarantees} (\cref{thm:unbiased,thm:contraction}). Anchored Energy Matching (AEM) adds $x_0$ as an extra particle. With an exact anchor the target is unbiased for every number of proposals $K\ge0$, its covariance is bounded by that of single-sample score matching, and $\pi$ becomes a fixed point. With an inexact buffer, the target equals the posterior mean after one i-SIR step from the buffer posterior, so its bias contracts the buffer's error by an explicit Doeblin factor. AEM passes the consistency test in every setting we tried, although the test has little power against it where the anchor dominates.
\item \textbf{A clear account of what consistency does and does not buy.} In full sampler training, AEM removes iDEM's within-mode distortion. It does not discover modes, and the contraction constant predicts why: where no proposal reaches a mode, the anchored target follows the buffer. iDEM explores better precisely because its target is biased.
\end{itemize}

In full sampler training on a GPU (3 seeds), iDEM's bias appears as a distorted within-mode distribution: on MW-32 its Gaussian coordinates have variance \val{mw32}{idem}{within_var}{m} instead of $1$ and its energy histogram is shifted. AEM removes this ($\val{mw32}{aem_ea}{within_var}{m}$--$\val{mw32}{aem_ea_refine}{within_var}{m}$) at $\ratio{evals_over_idem:mw32:aem_ea}$--$\ratio{evals_over_idem:mw32:aem_ea_refine}\times$ the energy evaluations.
We also report where consistency does not translate into better samplers. For mode discovery and mode weights, iDEM's biased target explores better and AEM needs an exploration phase to keep up. In full training on DW-4 and LJ-13, anchoring gives no significant gain at our budgets, and iDEM is better on LJ-13, even though AEM's target passes the consistency test there and iDEM's does not. Even with an exact buffer, AEM learns a slightly worse LJ-13 score at low noise, where iDEM's bias is small and the anchored target's variance is larger. Our message is therefore twofold: energy-matching targets should be checked for consistency, which is cheap, and a consistent target is necessary but not sufficient; exploration has to come from elsewhere.
